\chapter{Complexity — Uniform Tableau Circuit}
%%% generated by the blueprint pipeline from TCSlib.Complexity.CircuitComplexity.UniformTableauCircuit
\section{Overview}

Declarations of the Lean module \texttt{TCSlib.Complexity.CircuitComplexity.UniformTableauCircuit}.

\section{Declarations}

\begin{definition}[Copy gate]
\lean{BoolCircuit.copyGate}
\label{BoolCircuit.copyGate}
\leanok
\uses{BoolCircuit.DAGGate}
For each natural number $v$, define the copy gate at $v$ to be an AND gate whose sole
input is the vertex indexed by $v$.
\end{definition}

\begin{definition}[Shifted gate]
\lean{BoolCircuit.shiftGate}
\label{BoolCircuit.shiftGate}
\leanok
\uses{BoolCircuit.DAGGate}
Given a natural number $B$ and a gate $g$, the shifted gate has the same label as $g$
and has each argument index increased by $B$.
\end{definition}

\begin{theorem}[Copy gate evaluation]
\lean{BoolCircuit.copyGate_eval}
\label{BoolCircuit.copyGate_eval}
\leanok
\uses{BoolCircuit.DAGGate.eval, BoolCircuit.copyGate}
\difficulty{1}
For any natural number $v$ and Boolean list $\mathrm{vals}$, the gate that copies the
value at index $v$ evaluates to $\mathrm{vals}[v]$, taking the value to be false when
the list has no entry at that index.
\end{theorem}

\begin{theorem}[Shifted embedding of a circuit list]
\lean{BoolCircuit.embedAll_range'_eq}
\label{BoolCircuit.embedAll_range'_eq}
\leanok
\uses{BoolCircuit.DAGCircuit, BoolCircuit.DAGGate.dedupArgs, BoolCircuit.DAGGate.remap, BoolCircuit.embedAll, BoolCircuit.embedGates, BoolCircuit.embedMap, BoolCircuit.shiftGate}
\difficulty{6}
Let $m,B,L$ be natural numbers and let $D_1,\ldots,D_r$ be circuits with $m$ inputs.
Embedding this list using source vertices $B,B+1,\ldots,B+m-1$ and offset $B+L$ gives
the same gate list as embedding it using source vertices $0,1,\ldots,m-1$ and offset
$L$, then adding $B$ to every input vertex of every gate.
\end{theorem}

\begin{definition}[Maximum width over a finite family]
\lean{BoolCircuit.tabG}
\label{BoolCircuit.tabG}
\leanok
\uses{BoolCircuit.kindWidth}
Let $\kappa$ be a finite type, let $m,w\in\mathbb{N}$, and let $F$ assign to each
$k\in\kappa$ a function from Boolean lists to Boolean lists. Define $T$ to be the
maximum, over $k\in\kappa$, of the width value associated with $m,w$, and $F(k)$.
\end{definition}

\begin{definition}[Tabulation width parameter]
\lean{BoolCircuit.tabW}
\label{BoolCircuit.tabW}
\leanok
\uses{BoolCircuit.tabG}
For a finite type $\kappa$, natural numbers $m,w$, and a function $F$ assigning a
Boolean-list transformation to each element of $\kappa$, define the tabulation width
parameter as $m$ plus the maximum tabulation gate width for $m,w,F$, plus $w$.
\end{definition}

\begin{theorem}[Kind width bounded by the universal width]
\lean{BoolCircuit.kindWidth_le_tabG}
\label{BoolCircuit.kindWidth_le_tabG}
\leanok
\uses{BoolCircuit.kindWidth, BoolCircuit.tabG}
\difficulty{1}
Let $\kappa$ be a finite type, let $m,w\in\mathbb{N}$, and let $F$ assign to each
$\kappa$-parameter a function from Boolean lists to Boolean lists. For every
$\kappa_0\in\kappa$, the total size of the gadgets for the $w$ output coordinates of
$F(\kappa_0,\cdot)$ is at most the maximum of those sizes over all parameters in
$\kappa$.
\end{theorem}

\begin{definition}[Tabulation block base]
\lean{BoolCircuit.tabBase}
\label{BoolCircuit.tabBase}
\leanok
\uses{BoolCircuit.tabW}
Let $\kappa$ be a finite type, let $m,w,n\in\mathbb{N}$, let $F$ assign a Boolean list
to each pair consisting of an element of $\kappa$ and a Boolean list, and let $z$ be a
Boolean list. If $T$ denotes the total width reserved for one program instruction, the
base index of instruction $i\in\mathbb{N}$ is $n+2+|z|+T(i+1)$.
\end{definition}

\begin{definition}[Tabulation source index]
\lean{BoolCircuit.tabSrc}
\label{BoolCircuit.tabSrc}
\leanok
\uses{BoolCircuit.BitSrc, BoolCircuit.BitSrc.Valid, BoolCircuit.tabBase, BoolCircuit.tabG}
Let $\kappa$ be finite, let $m,w,n\in\mathbb{N}$, let
$F:\kappa\to\mathrm{List}(\mathrm{Bool})\to\mathrm{List}(\mathrm{Bool})$, and let $z$ be
a Boolean list. For instruction index $i$ and source $s$, assign an index as follows
when $s$ is valid for input bound $|z|+n$, block width $w$, and instruction index $i$:
an input $k$ is assigned $n+2+k$ if $k<|z|$, and $k-|z|$ otherwise; a constant bit $b$
is assigned $n+1$ if true and $n$ if false; and block bit $(i',j)$ is assigned the base
index of instruction $i'$ plus $m$, the gadget width, and $j$. If the source is invalid,
assign $n$.
\end{definition}

\begin{definition}[Fixed instruction gates]
\lean{BoolCircuit.tabFixed}
\label{BoolCircuit.tabFixed}
\leanok
\uses{BoolCircuit.DAGGate, BoolCircuit.constGate, BoolCircuit.copyGate, BoolCircuit.embedAll, BoolCircuit.embedOffset, BoolCircuit.kindGadgets, BoolCircuit.kindWidth, BoolCircuit.tabG}
Let $\kappa$ be finite, let $m,w,n\in\mathbb{N}$, let
$F:\kappa\to\mathrm{List}(\mathrm{Bool})\to\mathrm{List}(\mathrm{Bool})$, let $z$ be a
Boolean list, and choose $\kappa_0\in\kappa$. Form the gate list by embedding the
gadgets for $\kappa_0$ using the first $m$ source indices, then appending enough
true-constant gates to bring the gadget region to the maximum gadget width, and finally
appending one copy gate for each of the $w$ output positions, referring to that
position's embedded output.
\end{definition}

\begin{definition}[Padded source list]
\lean{BoolCircuit.padSrc}
\label{BoolCircuit.padSrc}
\leanok
\uses{BoolCircuit.BitSrc}
Let $\kappa$ be finite, let $m,w,n\in\mathbb{N}$, let
$F:\kappa\to\mathrm{List}(\mathrm{Bool})\to\mathrm{List}(\mathrm{Bool})$, let $z$ be a
Boolean list, and let $\mathrm{srcs}$ be a list of bit sources. Append $m$
false-constant sources to $\mathrm{srcs}$ and retain only the first $m$ entries.
\end{definition}

\begin{definition}[Gadget gates for one instruction]
\lean{BoolCircuit.tabInstrGates}
\label{BoolCircuit.tabInstrGates}
\leanok
\uses{BoolCircuit.DAGGate, BoolCircuit.GInstr, BoolCircuit.copyGate, BoolCircuit.padSrc, BoolCircuit.shiftGate, BoolCircuit.tabBase, BoolCircuit.tabFixed, BoolCircuit.tabSrc}
Let $\kappa$ be finite, let $m,w,n\in\mathbb{N}$, let
$F:\kappa\to\mathrm{List}(\mathrm{Bool})\to\mathrm{List}(\mathrm{Bool})$, and let $z$ be
a Boolean list. For instruction index $i$ and instruction $\mathrm{ins}$, take the
padded list of its source bits and replace each source by a copy gate referring to its
tabulation source index. Append the fixed gate list for the instruction's kind, shifting
every vertex index in that list by the base index for instruction $i$.
\end{definition}

\begin{definition}[Tabulation preamble gates]
\lean{BoolCircuit.tabPre}
\label{BoolCircuit.tabPre}
\leanok
\uses{BoolCircuit.DAGGate, BoolCircuit.constGate}
Let $\kappa$ be finite, let $m,w,n\in\mathbb{N}$, let
$F:\kappa\to\mathrm{List}(\mathrm{Bool})\to\mathrm{List}(\mathrm{Bool})$, let $z$ be a
Boolean list, and let $W\in\mathbb{N}$. The preamble consists, in order, of a
false-constant gate, a true-constant gate, one constant gate for each bit of $z$, and
$W$ false-constant gates.
\end{definition}

\begin{definition}[Gate list for a tabulated program]
\lean{BoolCircuit.tabProgGates}
\label{BoolCircuit.tabProgGates}
\leanok
\uses{BoolCircuit.DAGGate, BoolCircuit.GInstr, BoolCircuit.tabInstrGates}
Let $\kappa$ be finite, let $m,w,n\in\mathbb{N}$, let
$F:\kappa\to\mathrm{List}(\mathrm{Bool})\to\mathrm{List}(\mathrm{Bool})$, and let $z$ be
a Boolean list. For a program consisting of instructions
$\mathrm{ins}_0,\ldots,\mathrm{ins}_{r-1}$, concatenate in order the gate lists for
those instructions, using instruction index $i$ for $\mathrm{ins}_i$. The empty program
produces the empty gate list.
\end{definition}

\begin{theorem}[Length of the fixed table circuit]
\lean{BoolCircuit.length_tabFixed}
\label{BoolCircuit.length_tabFixed}
\leanok
\uses{BoolCircuit.kindWidth, BoolCircuit.kindWidth_le_tabG, BoolCircuit.length_embedAll, BoolCircuit.tabFixed, BoolCircuit.tabG}
\difficulty{3}
Let $\kappa$ be a finite type, let $m,w\in\mathbb{N}$, and let $F$ assign a Boolean list
to each element of $\kappa$ and each Boolean list. For every $\kappa_0\in\kappa$, the
fixed table circuit associated with $\kappa_0$ has length equal to the maximum gadget
width over $\kappa$, plus $w$.
\end{theorem}

\begin{theorem}[Length of the padded source list]
\lean{BoolCircuit.length_padSrc}
\label{BoolCircuit.length_padSrc}
\leanok
\uses{BoolCircuit.BitSrc, BoolCircuit.padSrc}
\difficulty{1}
Let $\kappa$ be a finite type, $m,w,n\in\mathbb{N}$, let
$F:\kappa\to\mathrm{List}(\mathrm{Bool})\to\mathrm{List}(\mathrm{Bool})$, let $z$ be a
Boolean list, and let $srcs$ be a list of bit sources. The padded source list has length
$m$.
\end{theorem}

\begin{theorem}[Length of an instruction gate list]
\lean{BoolCircuit.length_tabInstrGates}
\label{BoolCircuit.length_tabInstrGates}
\leanok
\uses{BoolCircuit.GInstr, BoolCircuit.length_tabFixed, BoolCircuit.tabInstrGates, BoolCircuit.tabW}
\difficulty{3}
Let $\kappa$ be a finite type, let $m,w,n$ be natural numbers, let $F$ assign to each
$a\in\kappa$ a function from Boolean lists to Boolean lists, and let $z$ be a Boolean
list. Let $T$ be the maximum over $a\in\kappa$ of the total size of the fixed circuit
representations of the $w$ output coordinates of $F(a,\cdot)$, counting one output gate
for each coordinate, and set $W=m+T+w$. For every natural number $i$ and every
instruction over $\kappa$, the gate list constructed for that instruction has length
$W$.
\end{theorem}

\begin{theorem}[Length of the generated gate list]
\lean{BoolCircuit.length_tabProgGates}
\label{BoolCircuit.length_tabProgGates}
\leanok
\uses{BoolCircuit.GInstr, BoolCircuit.length_tabInstrGates, BoolCircuit.tabProgGates, BoolCircuit.tabW}
\difficulty{4}
Let $\kappa$ be a finite type, let $m,w,n$ be natural numbers, let $F$ assign a list of
Boolean values to each element of $\kappa$ and each input list of Boolean values, and
let $z$ be a list of Boolean values. For every natural number $i$ and every finite list
of instructions whose kinds lie in $\kappa$, the list of gates generated from those
instructions starting at index $i$ has length $(m+\mathrm{tabG}(m,w,F)+w)$ times the
number of instructions.
\end{theorem}

\begin{theorem}[Length of the table prefix]
\lean{BoolCircuit.length_tabPre}
\label{BoolCircuit.length_tabPre}
\leanok
\uses{BoolCircuit.tabPre}
\difficulty{3}
Let $z$ be a list of Boolean values and let $W\in\mathbb{N}$. The table prefix formed
from two constant gates, one constant gate for each entry of $z$, and $W$ additional
false constant gates has length $2+|z|+W$.
\end{theorem}

\begin{theorem}[Source indices lie below their block base]
\lean{BoolCircuit.tabSrc_lt}
\label{BoolCircuit.tabSrc_lt}
\leanok
\uses{BoolCircuit.BitSrc, BoolCircuit.BitSrc.Valid, BoolCircuit.tabBase, BoolCircuit.tabSrc, BoolCircuit.tabW}
\difficulty{5}
Let $\kappa$ be a finite type, let $m,w,n$ be nonnegative integers, let $F$ assign a
list of Boolean values to each pair consisting of an element of $\kappa$ and a Boolean
list, and let $z$ be a Boolean list. For every stage $i$ and every source $s$—an input
position, a Boolean constant, or a pair giving an earlier block and bit position—the
vertex index assigned to $s$ by the tabulation encoding is strictly less than the base
vertex index for stage $i$.
\end{theorem}

\begin{theorem}[Acyclicity of a tabulated instruction gadget]
\lean{BoolCircuit.gatesAcyclic_tabInstrGates}
\label{BoolCircuit.gatesAcyclic_tabInstrGates}
\leanok
\uses{BoolCircuit.GInstr, BoolCircuit.GatesAcyclic, BoolCircuit.GatesAcyclic.append, BoolCircuit.copyGate, BoolCircuit.embedOffset_lt, BoolCircuit.gatesAcyclic_embedAll, BoolCircuit.kindGadgets, BoolCircuit.kindWidth, BoolCircuit.kindWidth_le_tabG, BoolCircuit.length_embedAll, BoolCircuit.length_padSrc, BoolCircuit.padSrc, BoolCircuit.shiftGate, BoolCircuit.tabBase, BoolCircuit.tabFixed, BoolCircuit.tabInstrGates, BoolCircuit.tabSrc, BoolCircuit.tabSrc_lt}
\difficulty{6}
Let $m,w,n\in\mathbb{N}$, let $\kappa$ be a finite type, let $F$ assign a list of
Boolean outputs to each instruction kind and input list, let $z$ be a list of Boolean
values, and let $i$ and $\mathrm{ins}$ be an instruction index and instruction. The gate
list constructed for this instruction is acyclic relative to its tabulated base offset.
\end{theorem}

\begin{theorem}[Acyclicity of tabulated program gates]
\lean{BoolCircuit.gatesAcyclic_tabProgGates}
\label{BoolCircuit.gatesAcyclic_tabProgGates}
\leanok
\uses{BoolCircuit.GInstr, BoolCircuit.GatesAcyclic, BoolCircuit.GatesAcyclic.append, BoolCircuit.GatesAcyclic.nil, BoolCircuit.gatesAcyclic_tabInstrGates, BoolCircuit.length_tabInstrGates, BoolCircuit.tabBase, BoolCircuit.tabProgGates, BoolCircuit.tabW}
\difficulty{4}
Let $m,w,n\in\mathbb{N}$, let $\kappa$ be a finite type, let $F$ assign a list of
Boolean outputs to each instruction kind and input list, and let $z$ be a finite list of
Boolean values. For every index $i$ and finite instruction list $\mathrm{prog}$, the
concatenated gate lists for the instructions in $\mathrm{prog}$ are acyclic relative to
the tabulated base offset at $i$.
\end{theorem}

\begin{theorem}[Membership in the constant-gate prefix]
\lean{BoolCircuit.mem_tabPre}
\label{BoolCircuit.mem_tabPre}
\leanok
\uses{BoolCircuit.DAGGate, BoolCircuit.constGate, BoolCircuit.tabPre}
\difficulty{4}
Let $z$ be a finite Boolean list and let $W$ be a natural number. Every gate in the list
formed by placing a false constant gate, a true constant gate, one constant gate for
each entry of $z$, and then $W$ false constant gates is itself a constant gate.
\end{theorem}

\begin{theorem}[Acyclicity of the tabulation prefix]
\lean{BoolCircuit.gatesAcyclic_tabPre}
\label{BoolCircuit.gatesAcyclic_tabPre}
\leanok
\uses{BoolCircuit.GatesAcyclic, BoolCircuit.constGate, BoolCircuit.mem_tabPre, BoolCircuit.tabPre}
\difficulty{4}
Let $n,W\in\mathbb{N}$ and let $z$ be a finite list of Boolean values. The list
consisting of gates for the two Boolean constants, one constant gate for each entry of
$z$, and $W$ additional false-constant gates is acyclic relative to the $n$ input
vertices.
\end{theorem}

\begin{definition}[Tabulated program circuit]
\lean{BoolCircuit.tabCircuit}
\label{BoolCircuit.tabCircuit}
\leanok
\uses{BoolCircuit.DAGCircuit, BoolCircuit.GInstr, BoolCircuit.GatesAcyclic.append, BoolCircuit.gatesAcyclic_tabPre, BoolCircuit.gatesAcyclic_tabProgGates, BoolCircuit.length_tabPre, BoolCircuit.length_tabProgGates, BoolCircuit.tabBase, BoolCircuit.tabG, BoolCircuit.tabPre, BoolCircuit.tabProgGates, BoolCircuit.tabW}
Let $\kappa$ be a finite type, let $m,w,n$ be natural numbers, let $F$ assign a
Boolean-list function to each element of $\kappa$, let $z$ be a Boolean list, and let
$\mathrm{prog}$ be a finite list of instructions, each labelled by an element of
$\kappa$ and carrying an ordered list of source bits. For natural numbers $o$ and $j$,
form the directed acyclic circuit whose gates are the preamble for $z$ and the
tabulation gates for $\mathrm{prog}$, in that order. Its output vertex is the vertex at
the base of instruction block $o$, shifted by $m$, the tabulation width, and $j$, if $o$
is a valid instruction index and $j<w$; otherwise its output vertex is $n$.
\end{definition}

\begin{theorem}[Size of the tabulated circuit]
\lean{BoolCircuit.tabCircuit_size}
\label{BoolCircuit.tabCircuit_size}
\leanok
\uses{BoolCircuit.DAGCircuit.size, BoolCircuit.GInstr, BoolCircuit.length_tabPre, BoolCircuit.length_tabProgGates, BoolCircuit.tabCircuit, BoolCircuit.tabW}
\difficulty{3}
Let $\kappa$ be a finite type, let $m,w,n$ be natural numbers, let $F$ assign a Boolean
list to each instruction kind and Boolean list, and let $z$ be a Boolean list. For every
program $P$ consisting of instructions with kinds in $\kappa$ and every
$o,j\in\mathbb{N}$, the size of the associated tabulated circuit is $n+2+|z|+W(|P|+1)$,
where $W$ is the tabulation width determined by $m,w,F$, and $|z|$ and $|P|$ denote list
lengths.
\end{theorem}

\begin{theorem}[Fan-in bound for the tabulated circuit]
\lean{BoolCircuit.tabCircuit_isFaninTwo}
\label{BoolCircuit.tabCircuit_isFaninTwo}
\leanok
\uses{BoolCircuit.DAGCircuit.IsFaninTwo, BoolCircuit.DAGGate, BoolCircuit.GInstr, BoolCircuit.constGate, BoolCircuit.copyGate, BoolCircuit.faninTwo_embedAll, BoolCircuit.gadget_isFaninTwo, BoolCircuit.kindGadgets, BoolCircuit.mem_tabPre, BoolCircuit.shiftGate, BoolCircuit.tabCircuit, BoolCircuit.tabFixed, BoolCircuit.tabInstrGates, BoolCircuit.tabProgGates}
\difficulty{6}
Let $\kappa$ be a finite set of instruction kinds, let $m,w,n$ be nonnegative integers,
let $F$ assign to each kind a function from Boolean lists to Boolean lists, and let $z$
be a Boolean list. For every finite program of instructions with kinds in $\kappa$, and
every pair of natural numbers $o,j$, the tabulated circuit is well formed: each gate has
distinct arguments, every NOT gate has exactly one argument, and every gate has at most
two arguments.
\end{theorem}

\begin{theorem}[Execution of copy gates]
\lean{BoolCircuit.runWith_copies}
\label{BoolCircuit.runWith_copies}
\leanok
\uses{BoolCircuit.DAGGate.eval, BoolCircuit.copyGate, BoolCircuit.copyGate_eval, BoolCircuit.runWith, BoolCircuit.runWith_cons}
\difficulty{5}
Let $V$ be a list of Boolean values and $v_s$ a list of natural-number indices, each
strictly less than the length of $V$. Executing, in order, the copy gates for the
indices in $v_s$ from $V$ yields $V$ followed by the values of $V$ at those indices, in
the order specified by $v_s$.
\end{theorem}

\begin{theorem}[Concatenation law for tabulation gates]
\lean{BoolCircuit.tabProgGates_append}
\label{BoolCircuit.tabProgGates_append}
\leanok
\uses{BoolCircuit.GInstr, BoolCircuit.tabProgGates}
\difficulty{4}
Let $\kappa$ be a finite set of instruction kinds, let $m,w,n$ be nonnegative integers,
let $F$ assign to each kind a function from Boolean lists to Boolean lists, and let $z$
be a Boolean list. For any two finite instruction lists $p,q$ and any starting position
$i$, the gates tabulated for their concatenation, starting at $i$, are the concatenation
of the gates for $p$ starting at $i$ and the gates for $q$ starting at $i+|p|$.
\end{theorem}

\begin{theorem}[Running the tabulation preamble]
\lean{BoolCircuit.runWith_tabPre}
\label{BoolCircuit.runWith_tabPre}
\leanok
\uses{BoolCircuit.DAGGate.eval, BoolCircuit.constGate, BoolCircuit.runWith, BoolCircuit.runWith_eval_constGates, BoolCircuit.tabPre}
\difficulty{3}
Let $z$ and $x$ be lists of Boolean values, and let $W$ be a nonnegative integer.
Running the tabulation preamble formed from the constant gates for false and true, the
constant gates for the entries of $z$, and $W$ further false constant gates, from
initial list $x$, yields $x$ followed by $[false,true]$, then $z$, then $W$ copies of
$false$.
\end{theorem}

\begin{definition}[Validity of a table configuration]
\lean{BoolCircuit.TabInv}
\label{BoolCircuit.TabInv}
\leanok
\uses{BoolCircuit.GInstr, BoolCircuit.progBlocks, BoolCircuit.tabBase, BoolCircuit.tabG}
Let $m,w,n$ be natural numbers, let $F$ assign a Boolean-list function to each element
of a finite type of instruction kinds, let $z$ be a Boolean list, let
$x:\text{Fin}(n)\to\{\mathrm{false},\mathrm{true}\}$, and let $P$ be a program. Write
$B_r$ for the table base offset at step $r$, $g$ for the gate-layout offset, and $C_r$
for the block computed by the first $r$ instructions of $P$ on input $z$ followed by the
bits of $x$. A list $V$ is a valid configuration at step $i$ if its length is the
prescribed table length at that step; its first $n$ entries encode $x$, entries $n$ and
$n+1$ are false and true, respectively, its next $|z|$ entries encode $z$, and for every
$r<i$ and $j<w$, the entry at position $B_r+m+g+j$ equals bit $j$ of $C_r$, with
out-of-range list accesses interpreted using the stated default values.
\end{definition}

\begin{theorem}[Evaluation of a tabulated source bit]
\lean{BoolCircuit.tabSrc_eval}
\label{BoolCircuit.tabSrc_eval}
\leanok
\uses{BoolCircuit.BitSrc, BoolCircuit.BitSrc.Valid, BoolCircuit.BitSrc.eval, BoolCircuit.GInstr, BoolCircuit.TabInv, BoolCircuit.progBlocks, BoolCircuit.tabSrc}
\difficulty{4}
Let $\kappa$ be a finite set of instruction kinds, let $m,w,n$ be nonnegative integers,
let $F$ assign to each kind a function from Boolean lists to Boolean lists, let $z$ be a
Boolean list, and let $x$ be an assignment of $n$ Boolean inputs. If a list of Boolean
vertex values satisfies the tabulation invariant at position $i$, then for every source
bit valid at position $i$ relative to the inputs formed by adjoining $z$ to the $n$
input bits and output width $w$, the value at the vertex assigned to that source bit
equals its value when evaluated from those inputs and the blocks produced by running the
whole program.
\end{theorem}

\begin{theorem}[Successor step for the tabulation invariant]
\lean{BoolCircuit.tabInv_succ}
\label{BoolCircuit.tabInv_succ}
\leanok
\uses{BoolCircuit.BitSrc.eval, BoolCircuit.DAGGate.eval, BoolCircuit.GInstr, BoolCircuit.ProgWF, BoolCircuit.TabInv, BoolCircuit.constGate, BoolCircuit.copyGate, BoolCircuit.embedAll, BoolCircuit.embedAll_range'_eq, BoolCircuit.embedOffset, BoolCircuit.embedOffset_lt, BoolCircuit.embedWidth, BoolCircuit.gadget_eval, BoolCircuit.kindGadgets, BoolCircuit.kindWidth, BoolCircuit.kindWidth_le_tabG, BoolCircuit.padSrc, BoolCircuit.progBlocks, BoolCircuit.progBlocks_getD, BoolCircuit.runWith, BoolCircuit.runWith_append, BoolCircuit.runWith_copies, BoolCircuit.runWith_embedAll_getD, BoolCircuit.runWith_eval_constGates, BoolCircuit.runWith_getD_of_lt, BoolCircuit.shiftGate, BoolCircuit.tabBase, BoolCircuit.tabFixed, BoolCircuit.tabG, BoolCircuit.tabInstrGates, BoolCircuit.tabSrc, BoolCircuit.tabSrc_eval, BoolCircuit.tabSrc_lt, BoolCircuit.tabW}
\difficulty{7}
Let $\kappa$ be a finite set of instruction kinds, let $m,w,n$ be nonnegative integers,
let $F$ assign to each kind a function from Boolean lists to Boolean lists, let $z$ be a
Boolean list, and let $x$ be an assignment of $n$ Boolean inputs. Suppose every source
in each program instruction is valid relative to the inputs formed by adjoining $z$ to
the $n$ input bits, with output width $w$, and every instruction has exactly $m$
sources. If $i$ is a position in the program and a list of Boolean vertex values
satisfies the tabulation invariant after the first $i$ instructions, then running the
gates for instruction $i$ on that list produces vertex values satisfying the invariant
after the first $i+1$ instructions.
\end{theorem}

\begin{theorem}[Evaluation of the tabulated circuit]
\lean{BoolCircuit.tabCircuit_eval}
\label{BoolCircuit.tabCircuit_eval}
\leanok
\uses{BoolCircuit.DAGCircuit.eval, BoolCircuit.DAGCircuit.values, BoolCircuit.DAGGate.eval, BoolCircuit.GInstr, BoolCircuit.ProgWF, BoolCircuit.TabInv, BoolCircuit.progBlocks, BoolCircuit.runWith, BoolCircuit.runWith_append, BoolCircuit.runWith_tabPre, BoolCircuit.tabBase, BoolCircuit.tabCircuit, BoolCircuit.tabInv_succ, BoolCircuit.tabPre, BoolCircuit.tabProgGates, BoolCircuit.tabProgGates_append, BoolCircuit.tabW}
\difficulty{5}
Let $\kappa$ be a finite set of instruction kinds, let $m,w,n$ be nonnegative integers,
let $F$ assign to each kind a function from Boolean lists to Boolean lists, and let $z$
be a Boolean list. Consider a finite program whose instructions each have a kind and an
ordered list of source bits; suppose every source in instruction $i$ is valid when the
available input bits are $z$ followed by $n$ input bits and the output width is $w$, and
every instruction has exactly $m$ sources. For any input assignment
$x:\mathrm{Fin}\,n\to\mathrm{Bool}$, program position $o$ within the program, and output
coordinate $j<w$, the tabulated circuit's value is the $j$th bit of the block produced
by instruction $o$ when the program is run on the concatenation of $z$ and $x$, with
missing list entries interpreted as false.
\end{theorem}
